\begin{proof}
\begin{enumerate}[label={$\langle1\rangle$\arabic*.},leftmargin=*]
\item 
\noindent\textbf{Step 1.} Put
  \(V_k=H^0(Y,B\otimes\cO_Y(-kD))\), viewed inside $H^0(Y,B)$, and note
  that $V_{k+1}\subseteq V_k$.  The descending chain stabilizes.  If $s$ is
  in its stable value, then at a point of $D$ choose a local defining
  nonzerodivisor $g_D$ for the Cartier divisor and a local frame of $B$.
  Every coefficient of the germ of $s$ is divisible by $g_D^k$ for every
  $k$.  The Krull-intersection theorem in the regular local ring
  $\cO_{Y,y}$ gives $\bigcap_{k\ge1}g_D^k\cO_{Y,y}=0$, so the germ is zero at
  every point of $D$.  Hence $s$ vanishes on a nonempty open set and, by
  connectedness, on $Y$.  Thus $V_k=0$ for all sufficiently large $k$.

  The chosen point $P_0$ satisfies
  $H^0(Y,B\otimes P_0)=0$.  The Poincar\'e bundle gives, locally on
$\mathcal P$, a family to which Grauert's theorem applies; the local jump
  loci $\{P':h^0(Y,B\otimes P')\ge1\}$ glue to a closed analytic subset of
  $\mathcal P$ (Grauert's theorem; see \cite{Grauert1960D,BS76push}).  It is proper because it misses $P_0$, and irreducibility of
$\mathcal P$ makes its complement dense open.  Applying the same stabilization
argument to $B^\vee\otimes\omega_Y$ and using the chosen point $Q_0$ gives
$H^0(Y,B^\vee\otimes\omega_Y\otimes Q_0)=0$; dualization stability puts
$Q_0^\vee$ in $\mathcal P$.  Serre duality then gives
$H^n(Y,B\otimes Q_0^\vee)=0$.  The corresponding top-degree jump locus is
therefore also proper with dense open complement.  Intersecting the two
dense open complements proves the claim.

\end{enumerate}
\end{proof}
